%!TEX root = paper.tex

\begin{abstract}
Python security analysis is largely source-centric: package scanners, dependency analyzers, malware detectors, and program-analysis pipelines typically inspect Python source files, metadata, installation scripts, and dependency relations.  However, CPython executes compiled code objects, routinely materializes bytecode as \texttt{.pyc} files, and can load bytecode through source-less modules, cached imports, bundled applications, custom loaders, and marshalled code objects.  This creates a practical visibility gap: behavior may be present in executable bytecode even when corresponding source is missing, stale, mismatched, or difficult to recover.

This paper presents an empirical study of Python bytecode as a first-class security artifact.  We study how bytecode appears in real software distributions, whether it is visible from source, how well existing tools recover source-level representations, and how CPython behaves when processing unusual bytecode inputs.  Our methodology combines package-scale artifact extraction, source/bytecode correspondence analysis, decompilation and recompilation checks, runtime robustness testing across CPython versions, and source-reachability analysis for abnormal behaviors.  The study is designed to separate ecosystem visibility risks from interpreter robustness findings and to distinguish bytecode-only behavior from behavior that ordinary source-level tooling can observe.

Our results show that Python bytecode deserves explicit treatment in package security workflows.  Bytecode artifacts occur in multiple distribution forms, source correspondence is not guaranteed, decompilation-based recovery is incomplete, and bytecode-level testing reaches runtime behaviors that are not always reproducible from ordinary source.  These findings motivate bytecode-aware package scanning, explicit source/bytecode mismatch checks, more robust bytecode validation boundaries, and security tooling that treats compiled Python artifacts as analyzable inputs rather than disposable caches.

\end{abstract}
